\chapter{La STFT como marco de Gabor}\label{cap:stft}

\enproceso{Análisis y síntesis; redundancia $N/H$; ventana dual; COLA y NOLA; la ISTFT como proyección de mínimos cuadrados (Griffin–Lim) y por qué una STFT modificada no es una STFT; consistencia; el principio de incertidumbre aplicado al espectrograma; la malla física $(t_m,f_k)$ de Fourierizer.}
